% GENERATED FILE. Edit canonical lessons, then run npm run generate:books.
% canonical-source-hash: fnv1a64:6df53d1a0e3e06a3
% canonical-lessons: BN-C511-bason, BN-C511-dekchi, BN-C511-patro, BN-C511-dhala, BN-C511-dobano

\chapter{Pots and Pouring}
\label{ch:bn-pots-and-pouring}
\hlchaptermodality{511}

\begin{chapteropening}
\textbf{By the end of this chapter:} \emph{I can talk about pots and pans, pouring and dipping.}

Complete the last lesson of chapter 511: I can talk about pots and pans, pouring and dipping.
\end{chapteropening}

\section[\texorpdfstring{\bn{বাসন} (bāsôn)}{বাসন (bāsôn)}]{\bn{বাসন} (bāsôn) — dishes; pots and pans}

\label{lesson:BN-C511-bason}



\begin{quote}
Before the new one: say the Bengali for to wring out, then the Bengali for to scour.
\end{quote}

\subsection*{You'll want to know: \bn{বাসন}}
\textbf{\bn{বাসন}} — \emph{bāsôn} — \textquotedblleft{}dishes; pots and pans\textquotedblright{}.

\begin{quote}
\bn{খাওয়ার} \bn{পরে} \bn{আমি} \bn{বাসন} \bn{ধুই।}
\end{quote}

\begin{quote}
\emph{Khāoyār pôre āmi bāsôn dhui.}
\end{quote}

\begin{quote}
After eating I wash the dishes.
\end{quote}

\begin{cousinweb}
\textbf{\bn{বাসন}} (\emph{bāsôn}) probably goes back to Sanskrit \emph{vāsana}, \textquotedblleft{}a box, a casket\textquotedblright{}; CDIAL gives no Bengali form there, but lists Prakrit \emph{vāsaṇa}, \textquotedblleft{}a vessel, a plate\textquotedblright{}, and Hindi \emph{bāsan}, \textquotedblleft{}a pot, a dish\textquotedblright{} (CDIAL 11599).
\end{cousinweb}

\begin{grammarlens}[title={What you've built}]
One more word for describing what you lived and felt.
\end{grammarlens}

\subsection*{Your turn}
\begin{itemize}
\raggedright
  \item \emph{Say it:} \emph{bāsôn}
  \item \emph{Say it:} \emph{bāsôn}, once more
  \item \emph{Say it:} \emph{Khāoyār pôre āmi bāsôn dhui.} — After eating I wash the dishes.
  \item \emph{Recall:} say the Bengali for to wring out, then the Bengali for to scour, then say \emph{bāsôn} again
\end{itemize}

\begin{tcolorbox}[breakable,colback=teal!4,colframe=teal!35!black,title={Before you move on}]
What does \bn{বাসন} mean? (\textquotedblleft{}Dishes; pots and pans\textquotedblright{}.) Say it once more.
\end{tcolorbox}
\section[\texorpdfstring{\bn{ডেকচি} (ḍekchi)}{ডেকচি (ḍekchi)}]{\bn{ডেকচি} (ḍekchi) — a wide cooking pot}

\label{lesson:BN-C511-dekchi}



\begin{quote}
Before the new one: say the Bengali for to scour, then the Bengali for dishes.
\end{quote}

\subsection*{You'll want to know: \bn{ডেকচি}}
\textbf{\bn{ডেকচি}} — \emph{ḍekchi} — \textquotedblleft{}a wide cooking pot\textquotedblright{}.

\begin{quote}
\bn{দিদিমা} \bn{ডেকচিতে} \bn{খিচুড়ি} \bn{রান্না} \bn{করলেন।}
\end{quote}

\begin{quote}
\emph{Didimā ḍekchite khichuṛi rānnā kôrlen.}
\end{quote}

\begin{quote}
My grandmother cooked khichuri in a big pot.
\end{quote}

\begin{cousinweb}
\textbf{\bn{ডেকচি}} (\emph{ḍekchi}) is usually traced to Persian \emph{degcha}, \textquotedblleft{}a small cauldron\textquotedblright{}, a smaller form of \emph{deg}, \textquotedblleft{}a cauldron\textquotedblright{}.
\end{cousinweb}

\begin{grammarlens}[title={What you've built}]
One more word for describing what you lived and felt.
\end{grammarlens}

\subsection*{Your turn}
\begin{itemize}
\raggedright
  \item \emph{Say it:} \emph{ḍekchi}
  \item \emph{Say it:} \emph{ḍekchi}, once more
  \item \emph{Say it:} \emph{Didimā ḍekchite khichuṛi rānnā kôrlen.} — My grandmother cooked khichuri in a big pot.
  \item \emph{Recall:} say the Bengali for to scour, then the Bengali for dishes, then say \emph{ḍekchi} again
\end{itemize}

\begin{tcolorbox}[breakable,colback=teal!4,colframe=teal!35!black,title={Before you move on}]
What does \bn{ডেকচি} mean? (\textquotedblleft{}A wide cooking pot\textquotedblright{}.) Say it once more.
\end{tcolorbox}
\section[\texorpdfstring{\bn{পাত্র} (pātrô)}{পাত্র (pātrô)}]{\bn{পাত্র} (pātrô) — a container; a vessel}

\label{lesson:BN-C511-patro}



\begin{quote}
Before the new one: say the Bengali for dishes, then the Bengali for a wide cooking pot.
\end{quote}

\subsection*{You'll want to know: \bn{পাত্র}}
\textbf{\bn{পাত্র}} — \emph{pātrô} — \textquotedblleft{}a container; a vessel\textquotedblright{}.

\begin{quote}
\bn{আমি} \bn{চালটা} \bn{একটা} \bn{বড়} \bn{পাত্রে} \bn{রাখলাম।}
\end{quote}

\begin{quote}
\emph{Āmi chālṭā ekṭā bôṛo pātre rākhlām.}
\end{quote}

\begin{quote}
I kept the rice in a big container.
\end{quote}

\begin{cousinweb}
\textbf{\bn{পাত্র}} (\emph{pātrô}) is a tatsama, Sanskrit \emph{pātra}, \textquotedblleft{}a drinking vessel, a dish\textquotedblright{} (Monier-Williams).
\end{cousinweb}

\begin{grammarlens}[title={What you've built}]
One more word for describing what you lived and felt.
\end{grammarlens}

\subsection*{Your turn}
\begin{itemize}
\raggedright
  \item \emph{Say it:} \emph{pātrô}
  \item \emph{Say it:} \emph{pātrô}, once more
  \item \emph{Say it:} \emph{Āmi chālṭā ekṭā bôṛo pātre rākhlām.} — I kept the rice in a big container.
  \item \emph{Recall:} say the Bengali for dishes, then the Bengali for a wide cooking pot, then say \emph{pātrô} again
\end{itemize}

\begin{tcolorbox}[breakable,colback=teal!4,colframe=teal!35!black,title={Before you move on}]
What does \bn{পাত্র} mean? (\textquotedblleft{}A container; a vessel\textquotedblright{}.) Say it once more.
\end{tcolorbox}
\section[\texorpdfstring{\bn{ঢালা} (ḍhālā)}{ঢালা (ḍhālā)}]{\bn{ঢালা} (ḍhālā) — to pour}

\label{lesson:BN-C511-dhala}



\begin{quote}
Before the new one: say the Bengali for a wide cooking pot, then the Bengali for a container.
\end{quote}

\subsection*{You'll want to know: \bn{ঢালা}}
\textbf{\bn{ঢালা}} — \emph{ḍhālā} — \textquotedblleft{}to pour\textquotedblright{}.

\begin{quote}
\bn{আমি} \bn{জগ} \bn{থেকে} \bn{গ্লাসে} \bn{জল} \bn{ঢাললাম।}
\end{quote}

\begin{quote}
\emph{Āmi jôg theke glāse jôl ḍhāllām.}
\end{quote}

\begin{quote}
I poured water from the jug into the glass.
\end{quote}

Past (he/she …): \bn{তিনি} \bn{কাপে} \bn{চা} \bn{ঢাললেন}, \emph{tini kāpe chā ḍhāllen}.

\begin{cousinweb}
\textbf{\bn{ঢালা}} (\emph{ḍhālā}) comes from a reconstructed \emph{ḍhalati}, \textquotedblleft{}tilts\textquotedblright{}; CDIAL lists Bengali \emph{ḍhālā}, \textquotedblleft{}to pour\textquotedblright{} (CDIAL 5581).
\end{cousinweb}

\begin{grammarlens}[title={What you've built}]
One more word for describing what you lived and felt.
\end{grammarlens}

\subsection*{Your turn}
\begin{itemize}
\raggedright
  \item \emph{Say it:} \emph{ḍhālā}
  \item \emph{Say it:} \emph{ḍhālā}, once more
  \item \emph{Say it:} \emph{Āmi jôg theke glāse jôl ḍhāllām.} — I poured water from the jug into the glass.
  \item \emph{Recall:} say the Bengali for a wide cooking pot, then the Bengali for a container, then say \emph{ḍhālā} again
\end{itemize}

\begin{tcolorbox}[breakable,colback=teal!4,colframe=teal!35!black,title={Before you move on}]
What does \bn{ঢালা} mean? (\textquotedblleft{}To pour\textquotedblright{}.) Say it once more.
\end{tcolorbox}
\section[\texorpdfstring{\bn{ডোবানো} (ḍobāno)}{ডোবানো (ḍobāno)}]{\bn{ডোবানো} (ḍobāno) — to dip; to sink (something)}

\label{lesson:BN-C511-dobano}



\begin{quote}
Before the new one: say the Bengali for a container, then the Bengali for to pour.
\end{quote}

\subsection*{You'll want to know: \bn{ডোবানো}}
\textbf{\bn{ডোবানো}} — \emph{ḍobāno} — \textquotedblleft{}to dip; to sink (something)\textquotedblright{}.

\begin{quote}
\bn{ছেলেটা} \bn{বিস্কুটটা} \bn{চায়ে} \bn{ডোবাল।}
\end{quote}

\begin{quote}
\emph{Chheleṭā biskuṭṭā chāye ḍobālô.}
\end{quote}

\begin{quote}
The boy dipped the biscuit in his tea.
\end{quote}

Past (he/she …): \bn{তিনি} \bn{রুটিটা} \bn{ডালে} \bn{ডোবালেন}, \emph{tini ruṭiṭā ḍāle ḍobālen}.

\begin{cousinweb}
\textbf{\bn{ডোবানো}} (\emph{ḍobāno}) is the causative partner of \bn{ডোবা}, \textquotedblleft{}to sink\textquotedblright{}, from a reconstructed \emph{ḍubb} (CDIAL 5561): to make something go under.
\end{cousinweb}

\begin{grammarlens}[title={What you've built}]
One more word for describing what you lived and felt.
\end{grammarlens}

\subsection*{Your turn}
\begin{itemize}
\raggedright
  \item \emph{Say it:} \emph{ḍobāno}
  \item \emph{Say it:} \emph{ḍobāno}, once more
  \item \emph{Say it:} \emph{Chheleṭā biskuṭṭā chāye ḍobālô.} — The boy dipped the biscuit in his tea.
  \item \emph{Recall:} say the Bengali for a container, then the Bengali for to pour, then say \emph{ḍobāno} again
\end{itemize}

\begin{tcolorbox}[breakable,colback=teal!4,colframe=teal!35!black,title={Before you move on}]
What does \bn{ডোবানো} mean? (\textquotedblleft{}To dip; to sink (something)\textquotedblright{}.) Say it once more.
\end{tcolorbox}
